%==============================================================================
\chapter{Design Implementation}
\label{ch:design}
%==============================================================================

This chapter details the realized design: the system architecture
(Section~\ref{sec:sysarch}), the domain model and lesson manifest
(Section~\ref{sec:domainmodel}), the backend implementation
(Section~\ref{sec:backendimpl}), the REST API (Section~\ref{sec:apidesign}),
the Android client (Section~\ref{sec:clientimpl}), the design patterns used
at the seams (Section~\ref{sec:patterns}), and the engineering
infrastructure --- hermetic testing, containerization, and CI/CD
(Section~\ref{sec:infra}).

%------------------------------------------------------------------------------
\section{System Architecture}
\label{sec:sysarch}
%------------------------------------------------------------------------------

Figure~\ref{fig:arch} shows the shape of the system: a relatively thin,
offline-first Android client talking over HTTPS to a stateless FastAPI
backend, which orchestrates hosted AI services through swappable adapters.

\begin{figure}[htbp]
\centering
\begin{tikzpicture}[
  font=\small,
  box/.style={draw, rounded corners, align=center, minimum height=8mm,
              inner sep=4pt, fill=blue!4},
  ext/.style={draw, rounded corners, align=center, minimum height=8mm,
              inner sep=4pt, fill=orange!8},
  grp/.style={draw, dashed, rounded corners, inner sep=8pt},
  >={Stealth}
]
% Client column
\node[box] (ui)   {UI / Screens\\(Compose)};
\node[box, below=5mm of ui]   (vm)  {ViewModels\\(StateFlow)};
\node[box, below=5mm of vm]   (uc)  {UseCases\\(domain)};
\node[box, below=5mm of uc, xshift=-14mm] (local) {Local Repo\\(Room + Files)};
\node[box, below=5mm of uc, xshift=16mm]  (remote){Remote Repo\\(Retrofit)};

\begin{scope}[on background layer]
\node[grp, fit=(ui)(vm)(uc)(local)(remote), label=above:{\textbf{Android client (offline-first)}}] (client) {};
\end{scope}

% Backend column
\node[box, right=38mm of vm] (rest) {REST API\\+ Job controller};
\node[box, below=5mm of rest] (facade) {Generation\\Facade};
\node[box, below=5mm of facade, xshift=-14mm] (graph) {LangGraph\\Agent};
\node[box, below=5mm of facade, xshift=18mm]  (cache) {Cache +\\Asset repos};

\begin{scope}[on background layer]
\node[grp, fit=(rest)(facade)(graph)(cache), label=above:{\textbf{FastAPI backend (stateless + shared cache)}}] (backend) {};
\end{scope}

% External column
\node[ext, right=26mm of rest] (llm) {Gemini\\LLM};
\node[ext, below=8mm of llm] (tts) {Gemini\\TTS};

\begin{scope}[on background layer]
\node[grp, fit=(llm)(tts), label=above:{\textbf{Hosted AI}}] (ai) {};
\end{scope}

% edges
\draw[->] (ui) -- (vm);
\draw[->] (vm) -- (uc);
\draw[->] (uc) -- (local);
\draw[->] (uc) -- (remote);
\draw[->] (remote) -- node[above, sloped, font=\footnotesize]{HTTPS} (rest);
\draw[->] (rest) -- (facade);
\draw[->] (facade) -- (graph);
\draw[->] (facade) -- (cache);
\draw[->] (graph) -- (llm);
\draw[->] (graph) -- (tts);
\end{tikzpicture}
\caption{System architecture. A thin offline-first Android client talks to a
stateless FastAPI backend over HTTPS. The backend orchestrates hosted AI
through swappable provider adapters and serves repeated topics from a shared
cache.}
\label{fig:arch}
\end{figure}

In deployment terms, the backend runs as a single container (Uvicorn +
FastAPI) whose generation work executes as asynchronous background tasks; the
lesson cache and asset store run, per configuration, either in memory, on the
filesystem, or against a database and object storage. The interfaces are
identical across these profiles, so promotion from the development profile to
a separate worker with cloud storage requires no call-site changes. The
production instance is deployed on a cloud virtual machine behind an
auto-HTTPS reverse proxy (Section~\ref{sec:infra}).

%------------------------------------------------------------------------------
\section{Domain Model and Lesson Manifest}
\label{sec:domainmodel}
%------------------------------------------------------------------------------

\subsection{Core Entities}

The backend's domain layer defines two central entities.

A \textbf{Lesson} carries an identifier, the original topic and its
normalized \texttt{topic\_key}, a display title, language and voice
parameters, the SVG document, the ordered list of segments, the total
duration, and a creation timestamp. A \textbf{Segment} carries its index, the
narration text, the list of SVG element identifiers it references, a
reference to its audio asset, and its measured duration in milliseconds.

Four invariants are enforced on these entities:

\begin{enumerate}
  \item \texttt{Lesson.svg} is well-formed XML with a \texttt{viewBox};
  \item every value in \texttt{Segment.svg\_element\_ids} exists as an
        \texttt{id} in \texttt{Lesson.svg};
  \item $\texttt{total\_duration\_ms} = \sum_i \texttt{duration\_ms}_i$; and
  \item segments are contiguous and ordered by index starting at 0.
\end{enumerate}

Invariants (1) and (2) are exactly what the deterministic gate of
Section~\ref{sec:gate} checks; (3) and (4) are established by the assembler.

\subsection{The Lesson Manifest}

The manifest is the wire and storage contract shared by the API and the
device (Listing~\ref{lst:manifest}). Asset URLs are downloadable; after
download the device rewrites them to local file paths, which is what makes
offline replay a non-event.

\begin{figure}[htbp]
\begin{lstlisting}[style=code,
  caption={Lesson manifest (abridged): the wire and storage contract.},
  label={lst:manifest}]
{
  "id": "les_8f3a...",
  "topic": "Photosynthesis",
  "title": "How Photosynthesis Works",
  "language": "en",
  "voice": "en-US-neutral",
  "total_duration_ms": 41200,
  "svg": { "asset_id": "ast_svg_01", "url": ".../assets/ast_svg_01" },
  "segments": [
    {
      "index": 0,
      "text": "Sunlight strikes the leaf and is captured by the plant.",
      "svg_element_ids": ["sun", "leaf"],
      "audio": { "asset_id": "ast_aud_00", "url": ".../assets/ast_aud_00" },
      "duration_ms": 4200
    },
    {
      "index": 1,
      "text": "Chlorophyll inside the chloroplasts absorbs that light energy.",
      "svg_element_ids": ["chloroplast", "chlorophyll"],
      "audio": { "asset_id": "ast_aud_01", "url": ".../assets/ast_aud_01" },
      "duration_ms": 3100
    }
  ]
}
\end{lstlisting}
\end{figure}

\subsection{Persistence Schemas}

On the backend, three stores back the service: \texttt{CACHED\_LESSON}
(keyed by \texttt{topic\_key}, holding the manifest and a hit counter),
\texttt{GENERATION\_JOB} (job id, topic key, status, progress, resulting
lesson id or error), and \texttt{ASSET} (asset id, kind
$\in\{\texttt{svg},\texttt{audio}\}$, content type, storage path). On the
device, Room stores \texttt{LESSON\_ENTITY} (metadata plus local SVG path)
and \texttt{SEGMENT\_ENTITY} (text, JSON-encoded identifier list, local
audio path, duration), while the binary assets live in app-private files.

%------------------------------------------------------------------------------
\section{Backend Implementation}
\label{sec:backendimpl}
%------------------------------------------------------------------------------

The backend is a Python FastAPI application of roughly two thousand lines,
organized by responsibility:

\begin{itemize}
  \item \texttt{app/api} --- routers and dependency wiring;
  \item \texttt{app/core} --- typed settings, logging, error envelope;
  \item \texttt{app/domain} --- entities, value objects, Pydantic schemas;
  \item \texttt{app/agents} --- the LangGraph graph and prompts;
  \item \texttt{app/providers} --- LLM and TTS adapters (mock and Gemini);
  \item \texttt{app/services} --- the generation facade and assembler; and
  \item \texttt{app/repositories} --- lesson cache and asset stores.
\end{itemize}

A single composition root wires these together, so swapping any
implementation happens in one place plus the typed settings object.

\subsection{Configuration over Hard-Coding}

Everything that varies between environments is read from environment
variables through one typed settings object, following the twelve-factor
configuration guideline~\cite{twelvefactor}: the choice of LLM and TTS
provider (mock or Gemini), the generation engine, the model identifiers, the
storage backend (memory or filesystem), and the loop budgets $K$ and $R$.
Model identifiers deserve emphasis: hosted model names change over time, so
they are configuration, not code. A small helper
(\texttt{python -m app.tools.list\_models}) lists the live models so the
pinned name can be confirmed; at the time of writing the backend is pinned to
a Flash-class Gemini text model (\texttt{gemini-2.5-flash}) and a Flash-class
Gemini speech model, both set in settings. The same codebase runs on mock or
real providers depending only on settings --- the property that keeps the
test suite fast and key-free.

\subsection{Provider Adapters}

Two narrow interfaces define what the pipeline needs. The \texttt{LLMProvider}
can plan concepts, draw an SVG, critique and refine an SVG, write narration
as structured data, and repair a drawing-plus-narration given an error list.
The \texttt{TTSProvider} turns text into an audio clip and reports its
duration. Each has two implementations:

\begin{itemize}
  \item \textbf{Mock providers} are deterministic and need no network ---
        ideal for tests and local development.
  \item \textbf{Gemini providers} call the hosted models through the
        \texttt{google-genai} SDK, using schema-constrained structured output
        for plans, narration, and repairs. The TTS adapter receives raw PCM
        audio, wraps it as a WAV file, and computes the exact duration from
        the byte count and sample rate --- the duration is a measured fact,
        not an estimate. Both adapters retry transient errors with
        exponential backoff.
\end{itemize}

A provider factory builds a coherent set of providers from settings
(abstract factory), and the generation facade
(\texttt{LessonGenerationFacade}) exposes one entry point that checks the
cache, runs the graph if needed, and persists the result.

\subsection{The Generation Graph in Code}

The pipeline of Section~\ref{sec:genmethod} is implemented as a LangGraph
\texttt{StateGraph} over a typed state object whose fields are set
progressively by the nodes: \texttt{topic} and \texttt{options} at entry;
\texttt{plan} by the planner; \texttt{svg} by the generator and refiner;
\texttt{segments} by the narration generator and repairer;
\texttt{validation} by the validator; \texttt{retries} by the repair loop;
\texttt{rendered} (audio plus durations) by the TTS renderer; \texttt{lesson}
by the assembler; and \texttt{error} by any node on hard failure. Conditional
edges route on the critique score and the validation verdict, with both loops
bounded by settings. Because the graph drives the provider \emph{interfaces},
the identical graph runs on mocks in tests and on Gemini in production.

%------------------------------------------------------------------------------
\section{API Design}
\label{sec:apidesign}
%------------------------------------------------------------------------------

The API is JSON over HTTPS under a versioned base path (\texttt{/api/v1}).
Generation is asynchronous because a cold lesson takes seconds, while cache
hits return immediately. Table~\ref{tab:api} lists the endpoints.

\begin{table}[htbp]
\caption{REST API endpoints}
\label{tab:api}
\centering
\small
\renewcommand{\arraystretch}{1.3}
\begin{tabular}{@{}llp{0.30\textwidth}p{0.24\textwidth}@{}}
\toprule
\textbf{Method} & \textbf{Path} & \textbf{Purpose} & \textbf{Success}\\
\midrule
POST & \texttt{/lessons} & Request a lesson for a topic & 200 (cache hit,
manifest) or 202 (job started)\\
GET & \texttt{/lessons/jobs/\{id\}} & Poll job status & 200 (status,
progress, stage)\\
GET & \texttt{/lessons/\{id\}} & Fetch a lesson manifest & 200 (manifest)\\
GET & \texttt{/assets/\{id\}} & Download an asset (SVG/audio) & 200 (bytes)\\
GET & \texttt{/healthz} & Liveness/readiness & 200\\
\bottomrule
\end{tabular}
\end{table}

A job response reports \texttt{status} $\in$ \{\texttt{queued},
\texttt{running}, \texttt{succeeded}, \texttt{failed}\}, a progress
percentage, and the current pipeline stage (\texttt{planning},
\texttt{svg}, \texttt{narration}, \texttt{validating},
\texttt{tts\_rendering}, \texttt{assembling}), which the client surfaces in
its progress UI. All errors share one envelope --- a machine-readable
\texttt{code} (for example \texttt{INVALID\_TOPIC},
\texttt{GENERATION\_FAILED}, \texttt{PROVIDER\_UNAVAILABLE}), a human
message, a \texttt{retryable} flag, and a correlation \texttt{request\_id}
--- so the client has exactly one error-handling path.
\texttt{POST /lessons} is idempotent per topic key, and concurrent identical
requests coalesce onto a single job (Section~\ref{sec:cachemethod}).

%------------------------------------------------------------------------------
\section{Android Client Implementation}
\label{sec:clientimpl}
%------------------------------------------------------------------------------

The client is a Kotlin application of roughly 3\,800 lines using Jetpack
Compose, structured by Clean Architecture~\cite{martin2017clean} into three
layers with dependencies pointing inward:

\begin{itemize}
  \item \textbf{Presentation} --- Compose screens (Topic, Player, History)
        and ViewModels that expose a single immutable \texttt{UiState} via
        \texttt{StateFlow}; the UI is a pure function of state.
  \item \textbf{Domain} --- pure Kotlin use cases (generate lesson, get
        history, play lesson, delete lesson), entities, and repository
        interfaces; no Android dependencies.
  \item \textbf{Data} --- a remote repository (Retrofit) implementing
        generation and download against the API, and a local repository
        (Room + file store) implementing persistence; both implement the
        domain's interfaces (dependency inversion).
\end{itemize}

Dependency injection is performed by a small manual composition root
(\texttt{AppContainer}); this was a deliberate deviation from the initial
Hilt plan, chosen for build reliability, and is swappable later since the
seams are the repository interfaces.

\subsection{Generation Flow on the Device}

The Topic screen submits the topic and then polls the job endpoint with
backoff, rendering the reported stage and progress. On success it fetches the
manifest and hands it to the player. Errors surface the envelope's message
with a retry affordance (FR-12).

\subsection{Synchronized Playback}

The player is the product's core (Figure~\ref{fig:player}). Two components
cooperate:

\begin{itemize}
  \item \textbf{Highlight engine.} The SVG is rendered in a WebView as an
        inline document together with a small app-owned script exposing
        \texttt{highlight(ids)} and \texttt{clearHighlight()}, which toggle a
        CSS class producing a glow-and-scale emphasis. The WebView renders
        \emph{only} this app-generated page; generated SVG is treated as
        untrusted markup and no content-supplied script is executed.
  \item \textbf{Segment clock.} ExoPlayer is loaded with one media item per
        segment clip, in order. The player's current-item index \emph{is} the
        synchronization state: on each transition to item $i$, the ViewModel
        highlights $E_i$ and updates the caption to the segment text.
\end{itemize}

Transport controls (previous, play/pause, next, replay) map onto playlist
operations, and the player follows the state machine
Idle $\to$ Loading $\to$ Ready $\to$ Playing $\rightleftharpoons$ Paused, with
Seeking, Completed (replay returns to Playing), and Error (retry returns to
Loading) states.

\begin{figure}[htbp]
\centering
\begin{tikzpicture}[
  font=\small,
  n/.style={draw, rounded corners, align=center, minimum height=8mm,
            inner sep=4pt, fill=blue!4},
  >={Stealth}
]
\node[n] (exo) {ExoPlayer\\playlist of $n$ clips};
\node[n, right=18mm of exo] (vmn) {PlayerViewModel\\(UiState)};
\node[n, right=18mm of vmn] (web) {WebView\\inline SVG + JS};
\node[n, below=8mm of vmn] (cap) {Caption\\(segment text)};

\draw[->] (exo) -- node[above,font=\footnotesize,align=center]
  {item $i$\\started} (vmn);
\draw[->] (vmn) -- node[above,font=\footnotesize,align=center]
  {highlight($E_i$)} (web);
\draw[->] (vmn) -- (cap);
\end{tikzpicture}
\caption{Synchronized playback. The audio player's current-item index drives
both the SVG highlight and the caption; no other timing state exists.}
\label{fig:player}
\end{figure}

\subsection{Offline Persistence and History}

``Save offline'' downloads the SVG and every audio clip into app-private
files and persists lesson and segment rows in Room; manifest asset URLs are
rewritten to local file URIs. The player is source-agnostic --- it receives a
loader function, so a saved lesson loads from local URIs through exactly the
same code path as a fresh one, and replays with zero network calls (NFR-2).
The History screen observes Room reactively (newest first); delete removes
both the rows and the files. History never leaves the device (NFR-7).

%------------------------------------------------------------------------------
\section{Design Patterns at the Seams}
\label{sec:patterns}
%------------------------------------------------------------------------------

Table~\ref{tab:patterns} maps the classical patterns~\cite{gamma1994} to
their locations and purposes in the system. The common motivation is
\emph{soft dependencies}: every external or replaceable concern sits behind
an interface chosen once, in one place.

\begin{table}[htbp]
\caption{Design patterns employed}
\label{tab:patterns}
\centering
\small
\renewcommand{\arraystretch}{1.3}
\begin{tabular}{@{}lp{0.34\textwidth}p{0.38\textwidth}@{}}
\toprule
\textbf{Pattern} & \textbf{Where} & \textbf{Why}\\
\midrule
Strategy & \texttt{LLMProvider}, \texttt{TTSProvider} & Swap AI vendors
without touching the graph\\
Abstract Factory & \texttt{ProviderFactory} & Build a coherent provider set
from configuration\\
Adapter & Gemini LLM/TTS classes & Wrap vendor SDKs behind the project's
interfaces\\
Builder & \texttt{LessonBuilder} & Assemble a valid \texttt{Lesson}
incrementally\\
Repository & Lesson cache, asset store, Android repositories & Hide storage
behind interfaces; memory/filesystem/DB profiles\\
Facade & \texttt{LessonGenerationFacade} & One entry point over cache check,
graph, and persistence\\
Pipeline & LangGraph nodes and conditional edges & Staged generation with
bounded loops\\
DTO & Pydantic schemas; manifest & Typed boundaries between layers and over
the wire\\
Dependency Injection & FastAPI wiring; \texttt{AppContainer} & Testability;
single composition root\\
Observer & \texttt{StateFlow} UiState; Room \texttt{Flow} & Reactive UI as a
pure function of state\\
\bottomrule
\end{tabular}
\end{table}

%------------------------------------------------------------------------------
\section{Hermetic Testing, Containerization, and CI/CD}
\label{sec:infra}
%------------------------------------------------------------------------------

\subsection{Hermetic, Key-Free Tests}

The backend suite comprises 48 tests covering the validator (both accepting
and rejecting), each graph node, the repair loop (recovery and
exhausted-budget failure), the API contract including the error envelope, the
job model, the repositories, and the Gemini adapters via a faked SDK client.
A session-level fixture disables \texttt{.env} loading and strips all
\texttt{TUTOR\_*} environment variables, so the suite is deterministic,
requires no secrets, and makes zero live API calls --- in any clone and in
CI.

\subsection{Containerization and Continuous Delivery}

The backend ships as a Docker image with a \texttt{/healthz} healthcheck and
a Compose profile for local orchestration. Four GitHub Actions workflows
automate the path from commit to artifact:

\begin{enumerate}
  \item \textbf{backend-ci} --- runs the hermetic test suite and builds the
        Docker image on every push;
  \item \textbf{backend-deploy} --- publishes the image to a container
        registry and deploys it over SSH to the cloud host, where an
        auto-HTTPS reverse proxy terminates TLS;
  \item \textbf{android-ci} --- assembles a debug APK as a build artifact;
        and
  \item \textbf{android-release} --- produces a signed release APK and
        attaches it to a GitHub Release, with a debug-key fallback when no
        signing secret is configured.
\end{enumerate}

Together with configuration-through-environment, this makes the entire system
reproducible: a fresh clone can run the full test suite with no keys, and a
tagged commit yields both a deployable backend image and an installable
client build with no manual steps.
